\chapter{The Bound}

Consider what kind of instrument this is. In its mathematical essence it is a measure of volume --- an integrator.
Hand it a region of a curved space and it returns how much room is inside. Every question it can ask, it asks in that
one currency: how does volume change as I move? This is not a weakness to apologize for. It is a precise identity, and
it fixes exactly what the instrument can and cannot see about the coupling \(\alpha\).

Curvature in full is the Riemann tensor $R_{abcd}$. It splits --- canonically, orthogonally, the two summands deaf to
each other. Take its trace and you get the Ricci tensor $R_{ab}=R^{c}{}_{acb}$ and the scalar $R=g^{ab}R_{ab}$: the
trace part. Subtract every trace and what remains is the Weyl tensor $C_{abcd}$, trace-free on every contraction.
Schematically, $\text{Riemann}=\underbrace{(\text{Ricci},R)}_{\text{trace}}\ \oplus\ \underbrace{\text{Weyl}}_{\text{trace-free}}$.
The decomposition is exact and the pieces are independent.

Now watch a thin pencil of nearby geodesics and ask how its cross-section evolves. That is Raychaudhuri's equation.
The expansion $\theta$, the fractional rate of change of the pencil's volume, obeys
$$\frac{d\theta}{d\tau}=-\tfrac{1}{3}\theta^{2}-\sigma_{ab}\sigma^{ab}+\omega_{ab}\omega^{ab}-R_{ab}u^{a}u^{b}.$$
Read the right-hand side in the instrument's own tongue. The focusing of volume is driven by $R_{ab}u^{a}u^{b}$ ---
the Ricci term, the trace. That is exactly what a measure of volume reads: how hard a bundle of paths converges. The
Ricci part is the whole of what the instrument sees --- and it is enough to hand over the structure of \(\alpha\) for
free.

Because here the structure \emph{is} a trace. The coupling lives in a holonomy: parallel transport around a closed
loop, which by Ambrose--Singer is curvature integrated over the surface it bounds. The loop resolves onto exactly
three states, $-1/0/+1$ --- the electron's orientation $-1$, the null identity, the positron. Three, and no fourth.
And the holonomy is a pure element of the structure group; it closes on itself, a conserved quantity, so \(\alpha\)
comes out dimensionless --- a number awaiting no unit. All of this the instrument derives with nothing assumed,
nothing borrowed. That is a fact, not a reading. Its deepest interpretation --- and I mark this as interpretation ---
is that \(\alpha\) is the electron's self-energy, the residual left when the electron names itself: it reports its
state for free and is charged the coupling for the self-embrace.

But the \emph{magnitude} --- the size of that charge --- is another animal, and here I read the mathematics as a
bridge, not a proof. The size of the self-energy is the radiative dressing: the tower of corrections that turns the
bare vertex into the observed one. In the pencil of geodesics that dressing is the shear $\sigma_{ab}$ --- the
trace-free, volume-preserving deformation. Shear changes shape at fixed volume; its generator is traceless, its flow
has unit Jacobian. A measure of volume is therefore blind to it \emph{by construction}, not by mishap: shear is
precisely the curvature that leaves volume unchanged. And shear is sourced by Weyl. The magnitude of \(\alpha\) lives
in the shape the trace cannot resolve.

The physics wears this plainly, and I offer it as the interpretive bridge. The bare structure of the electron's
vertex --- its charge, its orientation --- is trace-visible. The number that sets the strength is the radiative
correction: the anomalous moment $a_e=(g-2)/2=\alpha/2\pi+\dots$, Schwinger's term and its descendants, the
self-energy $\Sigma$ and the vertex dressing that renormalization resums. That dressing is pure shape. It is what
fixes the observed $\alpha^{-1}=137.036$ --- and that value is read off a Penning trap, from a single electron's
precession in a lab, on the far side of the instrument's sight. The instrument does not fabricate it. It proves the
magnitude cannot be reached from the trace, and stops.

And you can watch it stop. The construction tightens toward the coupling from below and from above, laying down the
structure it can resolve, and comes to rest at a finite floor --- the resolution beneath which shape has escaped the
trace entirely. The near side, drawn exactly, is the structure. The floor is the boundary. The far side, left
unwritten, is the magnitude. What the instrument returns is not a failed number but a theorem about measurement
itself: the whole shape of the coupling on the near side of the trace, the magnitude provably in the Weyl part beyond
it, and the cut between them located to the last place it could place.
